\section*{Блок 2. Центральная предельная теорема, Дельта-метод и сходимость по распределению}
\addcontentsline{toc}{section}{Блок 2. Центральная предельная теорема, Дельта-метод и сходимость по распределению}

\subsection*{1. Семантический анализ и математическая генерализация}
\addcontentsline{toc}{subsection}{1. Семантический анализ и математическая генерализация}

Глубокий анализ второго блока выявляет четкую иерархию задач: от базового понимания Центральной предельной теоремы (ЦПТ) до мощного аналитического инструмента — \textbf{Дельта-метода}. 

Закон больших чисел (ЗБЧ) устанавливает, что выборочное среднее $\overline{X} \toP \mu$. Однако ЗБЧ ничего не говорит о том, \textit{как именно} происходит эта сходимость. ЦПТ показывает, что флуктуации (отклонения) $\overline{X}$ от $\mu$ имеют порядок $1/\sqrt{n}$. Дельта-метод, в свою очередь, выступает мостом между асимптотикой самих средних и гладких функций от них, опираясь на разложение Тейлора первого порядка и теорему Слуцкого.

Опорная задача из КР и задачи 7, 8, 9, 10, 13 из семинаров реализуют один и тот же паттерн: требуется найти предельное распределение масштабированного отклонения $\sqrt{n}(g(\overline{X}) - C)$.

\subsection*{2. Фундаментальный инструментарий: Топология сходимостей}
\addcontentsline{toc}{subsection}{2. Фундаментальный инструментарий: Топология сходимостей}

Отношения между типами сходимости часто становятся объектом теоретических вопросов (задачи 1, 6, 11 и 12). Зафиксируем строгие математические границы их применимости:

\begin{itemize}
    \item \textbf{Связь с вероятностью:} Сходимость по распределению $\xi_n \toD \xi$ влечет сходимость по вероятности $\xi_n \toP \xi$ \textbf{тогда и только тогда}, когда предельная величина $\xi$ почти наверное есть константа $C$ (вырожденное распределение).
    \item \textbf{Отсутствие слабой компактности ЦПТ:} Последовательность $\eta_n = \sqrt{n}(\overline{X} - \mu)$ имеет предел по распределению, но \textbf{не имеет} предела по вероятности. Для любых $n \neq m$ разность $\eta_n - \eta_m$ не стремится к нулю по вероятности (последовательность не фундаментальна по Коши по вероятности).
    \item \textbf{Предельный переход в многомерных пространствах:} Из сходимости вектора по распределению следует сходимость компонент, но \textbf{не наоборот}. Из $\xi_n \toD \xi$ и $\eta_n \toD \eta$ в общем случае \textbf{не следует} $(\xi_n, \eta_n) \toD (\xi, \eta)$. Для совместной сходимости требуется либо независимость компонент при каждом $n$, либо применение многомерной ЦПТ к исходному случайному вектору.
\end{itemize}

\subsection*{3. Универсальный алгоритм решения (Многомерный Дельта-метод)}
\addcontentsline{toc}{subsection}{3. Универсальный алгоритм решения (Многомерный Дельта-метод)}



Разберем общий случай векторного Дельта-метода, покрывающий конструкции вида $T_n = Z_n / Y_n^k$, где $Z_n, Y_n$ — выборочные средние функций от выборки (задачи 10b, 13).

\begin{MethodBox}[Алгоритм применения векторного Дельта-метода]
Пусть требуется найти предел по распределению для $\sqrt{n} \left( g(Y_n, Z_n) - C \right)$.

\begin{enumerate}[label=\textbf{\arabic*.}]
    \item \textbf{Определение базиса.} Введите базовый случайный вектор $\vec{W}_i = (W_{1i}, W_{2i})^T$ для одного элемента выборки (например, $W_{1i} = |X_i|$, $W_{2i} = X_i^2$). Найдите вектор математических ожиданий $\vec{\mu} = \E[\vec{W}_1]$.
    
    \item \textbf{Проверка константы.} Убедитесь, что вычитаемая константа $C$ строго равна $g(\vec{\mu})$. Если это не так — предел бесконечен или равен нулю.
    
    \item \textbf{Матрица ковариаций.} Вычислите матрицу $\Sigma = \Cov(\vec{W}_1)$. На диагонали располагаются $\Var(W_{11})$ и $\Var(W_{21})$, вне диагонали — $\Cov(W_{11}, W_{21})$.
    
    \item \textbf{Вектор-градиент (Матрица Якоби).} Вычислите вектор частных производных функции $g$ в точке математического ожидания:
    \[
        \nabla g(\vec{\mu}) = \left( \frac{\partial g}{\partial y}(\vec{\mu}), \frac{\partial g}{\partial z}(\vec{\mu}) \right)^T
    \]
    
    \item \textbf{Синтез.} Асимптотическая дисперсия является квадратичной формой:
    \[
        \sigma^2_{as} = \nabla g(\vec{\mu})^T \cdot \Sigma \cdot \nabla g(\vec{\mu})
    \]
    Итоговый ответ: последовательность сходится по распределению к $\N(0, \sigma^2_{as})$.
\end{enumerate}
\vspace{0.5em}
\textbf{Аналитическое следствие для $g(y, z) = z / y^k$:} Градиент в точке $(\mu_y, \mu_z)$ имеет вид $\left( -k \frac{\mu_z}{\mu_y^{k+1}}, \frac{1}{\mu_y^k} \right)^T$. Итоговая дисперсия:
\[
    \sigma^2_{as} = \frac{k^2 \mu_z^2}{\mu_y^{2k+2}} \Var(W_y) - \frac{2k \mu_z}{\mu_y^{2k+1}} \Cov(W_y, W_z) + \frac{1}{\mu_y^{2k}} \Var(W_z)
\]
\end{MethodBox}

\subsection*{4. Эвристики вычислений и методы верификации}
\addcontentsline{toc}{subsection}{4. Эвристики вычислений и методы верификации}

Рассмотрим задачу №2 из КР: дано $X_i \sim Exp(\theta)$. Найти асимптотику $\sqrt{n}\left(\frac{1}{\overline{X}^2} - \theta^2\right)$.


\textbf{Решение:} Базовая ЦПТ: $$\sqrt{n}(\overline{X} - 1/\theta) \toD \N(0, 1/\theta^2)$$ 
\begin{enumerate}
    \item Функция $g(t) = t^{-2}$, $g'(t) = -2t^{-3}$.
    \item Градиент в точке $\mu = 1/\theta$: $g'(1/\theta) = -2\theta^3$.
    \item Асимптотическая дисперсия: $\sigma^2_{as} = (-2\theta^3)^2 \cdot (1/\theta^2) = 4\theta^4$.
    \item Ответ: $\N(0, 4\theta^4)$.
\end{enumerate}

Вычисление недиагональных элементов матрицы $\Sigma$ часто приводит к ошибкам интегрирования. Следующая лемма радикально упрощает процесс.

\begin{HackBox}[Симметрия абсолютных моментов (к задачам 10b, 13)]
Для нахождения $\Cov(|X|, X^2)$, если $X \sim \N(0, \sigma^2)$ или $X$ имеет распределение Лапласа, используйте симметрию плотности (все нечетные моменты равны нулю):
\[
    \Cov(|X|, X^2) = \E[|X| \cdot X^2] - \E|X| \cdot \E[X^2] = \E[|X|^3] - \E|X|\E[X^2]
\]
Раскрытие модуля не требуется! Выражение $\E[|X|^3]$ является абсолютным моментом третьего порядка. Для нормального распределения $\N(0, \sigma^2)$ справедлива формула:
\[
    \E|X|^p = \sigma^p \frac{2^{p/2} \Gamma\left(\frac{p+1}{2}\right)}{\sqrt{\pi}}
\]
\textbf{Готовые значения для экзамена:} $\E|X| = \sigma \sqrt{\frac{2}{\pi}}$, \quad $\E|X|^3 = 2\sigma^3 \sqrt{\frac{2}{\pi}}$.
Подстановка этих значений сводит вычисление ковариации к одной алгебраической строке.
\end{HackBox}

\begin{HackBox}[Sanity Checks: Методы быстрой самопроверки]
Перед записью финального ответа проведите аналитические тесты:
\begin{enumerate}[label=\textbf{\arabic*.}]
    \item \textbf{Анализ размерностей.} Если $X$ измеряется в [метрах], то $\Var(X)$ имеет размерность [метры$^2$]. Для функции $g(X) = 1/X^2$ дисперсия $\sigma^2_{as}$ обязана иметь размерность $[g(X)]^2$, то есть [метры$^{-4}$]. Нарушение размерности однозначно указывает на ошибку дифференцирования.
    \item \textbf{Строгая положительность.} Появление $\sigma^2_{as} < 0$ при расчете квадратичной формы $\nabla g^T \Sigma \nabla g$ указывает на ошибку в вычислении ковариации $\Cov(W_y, W_z)$ или в знаках матрицы Якоби.
\end{enumerate}
\end{HackBox}

\begin{HackBox}[Аналитический Cheat-Sheet: Моменты базовых распределений]
При решении задач асимптотического анализа, расчете ковариационных матриц и применении Дельта-метода критически важно оперировать не только математическим ожиданием и дисперсией, но и \textbf{вторым начальным моментом} $\E[X^2]$. Фундаментальное тождество для самопроверки: $\E[X^2] = \Var(X) + (\E X)^2$.

\vspace{0.5em}
\textbf{\textsf{1. Дискретные распределения}}
\vspace{0.2em}

\renewcommand{\arraystretch}{1.5}
\begin{tabular}{p{0.25\textwidth} p{0.2\textwidth} p{0.2\textwidth} p{0.25\textwidth}}
    \textbf{Распределение} & $\E X$ & $\Var(X)$ & $\E[X^2]$ \\
    \hline
    \textbf{Бернулли} $\Be(p)$ \newline $X \in \{0, 1\}$ 
    & $p$ 
    & $p(1-p)$ 
    & $p$ \\

    \textbf{Биномиальное} $\Bi(n, p)$ 
    & $np$ 
    & $np(1-p)$ 
    & $np(1-p) + n^2p^2$ \\

    \textbf{Пуассона} $\Po(\lambda)$ 
    & $\lambda$ 
    & $\lambda$ 
    & $\lambda(\lambda + 1)$ \\

    \textbf{Геометрическое} \newline $P(X=k) = p(1-p)^{k-1}$
    & $\frac{1}{p}$ 
    & $\frac{1-p}{p^2}$ 
    & $\frac{2-p}{p^2}$ \\
\end{tabular}

\vspace{1em}
\textbf{\textsf{2. Абсолютно непрерывные распределения}}
\vspace{0.2em}

\begin{tabular}{p{0.25\textwidth} p{0.2\textwidth} p{0.2\textwidth} p{0.25\textwidth}}
    \textbf{Распределение} & $\E X$ & $\Var(X)$ & $\E[X^2]$ \\
    \hline
    \textbf{Равномерное} $U(a, b)$ 
    & $\frac{a+b}{2}$ 
    & $\frac{(b-a)^2}{12}$ 
    & $\frac{a^2 + ab + b^2}{3}$ \\

    \textbf{Нормальное} $\N(\mu, \sigma^2)$ 
    & $\mu$ 
    & $\sigma^2$ 
    & $\mu^2 + \sigma^2$ \\

    \textbf{Экспоненциальное} \newline $p(x) = \lambda e^{-\lambda x}$ 
    & $\frac{1}{\lambda}$ 
    & $\frac{1}{\lambda^2}$ 
    & $\frac{2}{\lambda^2}$ \\

    \textbf{Гамма} $\Gamma(\alpha, \lambda)$ \newline $p(x) \propto x^{\alpha-1} e^{-\lambda x}$
    & $\frac{\alpha}{\lambda}$ 
    & $\frac{\alpha}{\lambda^2}$ 
    & $\frac{\alpha(\alpha+1)}{\lambda^2}$ \\

    \textbf{Лапласа} \newline $p(x) = \frac{1}{2b}e^{-|x-\mu|/b}$
    & $\mu$ 
    & $2b^2$ 
    & $\mu^2 + 2b^2$ \\
\end{tabular}

\vspace{1em}
\textbf{\textsf{3. Продвинутые моменты для $\N(0, \sigma^2)$ и распределения Коши}}
\vspace{0.2em}

Для стандартных симметричных распределений с нулевым средним ($\mu = 0$) полезно знать старшие и абсолютные моменты (необходимы для вычисления ковариаций модулей и квадратов):
\begin{itemize}
    \item \textbf{Старшие моменты $\N(0, \sigma^2)$:} $\E[X^3] = 0$ (в силу симметрии), $\E[X^4] = 3\sigma^4$.
    \item \textbf{Абсолютные моменты $\N(0, \sigma^2)$:} $\E|X| = \sigma \sqrt{\frac{2}{\pi}}$, \quad $\E|X|^3 = 2\sigma^3 \sqrt{\frac{2}{\pi}}$.
    \item \textbf{Распределение Коши $K(x_0, \gamma)$:} Математическое ожидание \textbf{не существует} (интеграл Лебега расходится). Дисперсия и старшие моменты равны бесконечности. Применение ЗБЧ и ЦПТ к выборке из распределения Коши строго запрещено. Выборочное среднее распределения Коши имеет то же самое распределение Коши, а не нормальное.
\end{itemize}
\end{HackBox}

\subsection*{5. Опасные места и аномалии сходимости}
\addcontentsline{toc}{subsection}{5. Опасные места и аномалии сходимости}

Классический Дельта-метод предполагает доминирование линейного члена разложения Тейлора. 

\begin{WarningBox}[Вырожденный градиент и Теорема о $\chi^2$ (Задача 10a)]
Если для функции $g(t)$ выполняется условие $g'(\mu) = 0$, но $g''(\mu) \neq 0$, классический Дельта-метод дает вырожденную дисперсию $\sigma^2_{as} = 0$, что означает лишь $\toP 0$. Множитель $\sqrt{n}$ становится недостаточным. 

Требуется переход к Дельта-методу второго порядка. Разложение Тейлора домножается не на $\sqrt{n}$, а на $n$:
\[
    n \left( g(\overline{X}_n) - g(\mu) \right) \approx \frac{1}{2} g''(\mu) \left[ \sqrt{n}(\overline{X}_n - \mu) \right]^2
\]
Поскольку $\sqrt{n}(\overline{X}_n - \mu) \toD \N(0, \sigma^2)$, квадрат этой величины сходится к $\sigma^2 \chi^2_1$ (хи-квадрат с одной степенью свободы). Итоговое предельное распределение:
\[
    n \left( g(\overline{X}_n) - g(\mu) \right) \toD \frac{\sigma^2 g''(\mu)}{2} \chi^2_1
\]
\textit{Пример:} Для $g(y) = y \sin y$, если точка матожидания $\mu = 0$, то $g'(0) = 0$, $g''(0) = 2$. Предельное распределение не будет нормальным — оно будет принадлежать семейству гамма-распределений (пропорционально $\chi^2_1$). Внимательно проверяйте условие $g'(\mu) \neq 0$.
\end{WarningBox}
